% problem_id: S046-thm:h-set-refrel-in-paths-sets
% source_id: S046 (HoTT Book)
% source_file: hlevels.tex
% statement_locator: hlevels.tex lines 265--269
% proof_locator: hlevels.tex lines 271--313
% env: thm
% pairing: tight (verified: the next environment after the statement is its proof)
% status: complete (statement + author proof verbatim + reason + locators)
%
\begin{thm}\label{thm:h-set-refrel-in-paths-sets}
  \index{relation!reflexive}%
  Suppose $R$ is a reflexive\index{reflexivity!of a relation} mere relation on a type $X$ implying identity.
  Then $X$ is a set, and $R(x,y)$ is equivalent to $\id[X]{x}{y}$ for all $x,y:X$.
\end{thm}

\begin{proof}
  Let $\rho : \prd{x:X} R(x,x)$ witness reflexivity of $R$, and let \narrowequation{f : \prd{x,y:X} R(x,y) \to (\id[X]{x}{y})} be a witness that $R$
implies identity.
  Note first that the two statements in the theorem are equivalent.
  For on one hand, if $X$ is a set, then $\id[X]xy$ is a mere proposition, and since it is logically equivalent to the mere proposition
$R(x,y)$ by hypothesis, it must also be equivalent to it.
  On the other hand, if $\id[X]xy$ is equivalent to $R(x,y)$, then like the latter it is a mere proposition for all $x,y:X$, and hence $X$
is a set.

  We give two proofs of this theorem.
  The first shows directly that $X$ is a set; the second shows directly that $R(x,y)\eqvsym (x=y)$.

  \emph{First proof:} we show that $X$ is a set.
  The idea is the same as that of \cref{thm:prop-set}: the function $f$ must be continuous in its arguments $x$ and $y$.
  However, it is slightly more notationally complicated because we have to deal with the additional argument of type $R(x,y)$.

  Firstly, for any $x:X$ and $p:\id[X]xx$, consider $\apdfunc{f(x)}(p)$.
  This is a dependent path from $f(x,x)$ to itself.
  Since $f(x,x)$ is still a function $R(x,x) \to (\id[X]xx)$, by \cref{thm:dpath-arrow} this yields for any $r:R(x,x)$ a path
  \[\trans{p}{f(x,x,r)} = f(x,x,\trans{p}r).
  \]
  On the left-hand side, we have transport in an identity type, which is concatenation.
  And on the right-hand side, we have $\trans{p}r = r$, since both lie in the mere proposition $R(x,x)$.
  Thus, substituting $r\defeq \rho(x)$, we obtain
  \[ f(x,x,\rho(x)) \ct p = f(x,x,\rho(x)). \]
  By cancellation, $p=\refl{x}$.
  So $X$ satisfies Axiom K, and hence is a set.

  \emph{Second proof:} we show that each $f(x,y) : R(x,y) \to \id[X]{x}{y}$ is an equivalence.
  By \cref{thm:total-fiber-equiv}, it suffices to show that $f$ induces an equivalence of total spaces:
  \begin{equation*}
    \eqv{\Parens{\sm{y:X}R(x,y)}}{\Parens{\sm{y:X}\id[X]{x}{y}}}.
  \end{equation*}
  By \cref{thm:contr-paths}, the type on the right is contractible, so it
  suffices to show that the type on the left is contractible. As the center of
  contraction we take the pair $\pairr{x,\rho(x)}$.  It remains to show, for
  every ${y:X}$ and every ${H:R(x,y)}$ that
  \begin{equation*}
    \id{\pairr{x,\rho(x)}}{\pairr{y,H}}.
  \end{equation*}
  But since $R(x,y)$ is a mere proposition, by \cref{thm:path-sigma} it suffices to show that
  $\id[X]{x}{y}$, which we get from $f(H)$.
\end{proof}
